\section{Similar Tasks}
There are numerous similar tasks in the literature. That is, applying machine-learning to the classification of security-sensitive methods. Here we will detail the more relevant papers in the literature.

\subsection{Swan}

SWAN is an automated classifier of security-sensitive methods in Java using machine-learning \cite{swanassist}. SWAN classifies code as a source, sink, sanitizer, authentication-method, or none --- which is identical to our task, with the inclusion of the \emph{authenitication-method} security-sensitive method type. SWAN is also able to further classify some of these methods according to their specific CWE (Common Weakness Enumeration), however this aspect of SWAN is not relevant to this thesis. \\

 \noindent
 SWAN employs a Support Vector Machine (SVM) approach on 206 manually defined binary features (i.e., does the method have this feature --- true or false?) to perform its classification. That is, the authors have gone through a dataset of examples of security-sensitive-methods and manually identified what aspects of a code is important for classifying it as a security sensitive method. This is decidedly different from our proposed implementation of fine-tuning a pre-trained bimodal Transformer model. \\

\noindent
For example, one of the 206 features SWAN classifies on is whether or not the class name the method belongs to contains the keyword `SQL' (the reason for this being that methods interacting with an SQL database are more likely to be a \emph{sink}). While this may seem like a simple feature, which it is, when combined with 205 additional features a high accuracy can be reached on a relatively small training dataset.\\

\noindent
SWAN's training dataset is publicly available and open-source \cite{swan_github}. Labelled datasets of security-sensitive methods are hard to come by, and thankfully Swan has quite an extensive list. For this thesis, we will be including a subset of this dataset (we remove \emph{authentication-method} examples, add \emph{none} examples which are not included in SWAN's public dataset, and remove unsuitable examples according to our own filtering rules) for our own training and evaluation purposes. Furthermore, SWAN does not provide its testing or validation dataset so we will split up their training set to get testing and validation examples.\\

\noindent
This dataset does not provide the code and Javadoc associated with each method, so additional processing of this data is required. We will come back to this later in the thesis. \\

\noindent
The \emph{SWAN} project also includes SWAN-assist, which is an Intellij plugin that allows users to classify security-relevant methods while working on a project. SWAN includes this additional information to improve the machine-learning classifier for a particular project. Of course, this aspect is not relevant to our project but it does inflate the accuracy of the model when compared to a fully-automated process such as ours. \\

\noindent
SWAN-base (i.e., SWAN without additional information from users using SWAN-assist) achieves an accuracy of $0.826$ on classification of security-sensitive methods \cite{swanassist}. While it would be good for our model to achieve an accuracy higher than this, it should be noted that it is not a fair comparison. For one, we will be operating on a subset of SWAN's base training-set out of necessity, and will be classifying on four labels, not five. Naturally, classifying on fewer labels is likely to inflate the accuracy of our model compared to SWAN-base, whereas operating on a subset on the dataset is likely to reduce the accuracy of our model relatively.

\subsection{SuSi}

